\section*{第 \ref{ch:g7:sound-production} 章 --- 声音：产生与传播}

\begin{solution}{exo:g7:sound-production:1}
声音所经过的物质——空气、水、木材，
钢铁，任何由分子聚集而成的物质。不是介质的：真空。
\end{solution}

\begin{solution}{exo:g7:sound-production:2}
接力在进行：一连串挤压脉冲，像接力棒一样从分子
一排传到下一排。每个分子只在自己的平衡位置附近
振动——信息在走，传信的分子却留在原地。
\end{solution}

\begin{solution}{exo:g7:sound-production:3}
没有分子就没有接力：抽掉空气的玻璃罩里，
钟的振动壁只能推着虚空，因此不会形成声音。太空
战斗之所以无声，原因相同——恒星之间的爆炸
没有分子传递它们的轰鸣。
\end{solution}

\begin{solution}{exo:g7:sound-production:4}
空气中约为 \qty{340}{m/s}；水中约为 \qty{1500}{m/s}；钢铁
中约为 \qty{5000}{m/s}。分子排列越紧密，
分子飞散、彼此接触、紧紧锁定，推力传递得越快。
\end{solution}

\begin{solution}{exo:g7:sound-production:5}
钢铁中分子紧密排列，接力几乎是手递手进行，
速率约为空气中的十五倍，因此沿钢轨传来的信息
远早于空气中传来的隆隆声。
\end{solution}

\begin{solution}{exo:g7:sound-production:6}
与声音相比，光穿过任何场地几乎可以看作瞬间到达：
看到的击打动作标记了真正的开始时刻，精度远高于
手动秒表所能达到的精度。
\end{solution}

\begin{solution}{exo:g7:sound-production:7}
$340 \times 1 = \qty{340}{m}$ 为往返路程，墙壁距离
\qty{170}{m}。若答 \qty{340}{m}，就忘了这一秒既包括去程，
\emph{又}包括回程。
\end{solution}

\begin{solution}{exo:g7:sound-production:8}
在水中：$1500 \times 2 = \qty{3000}{m}$ 为往返路程，深度
\qty{1500}{m}。错误在于把空气中的 \qty{340}{m/s} 用作
水下传播的速率。
\end{solution}

\begin{solution}{exo:g7:sound-production:9}
$d = 340 \times 2.5 = \qty{850}{m}$。烟花的闪光给出了一个确定的起跑信号；
而看不见的喷气机却没有这样的信号——无从确定
发声的时刻——更糟的是，听到的轰鸣声是喷气机在几秒前发出的，
耳朵指向的是喷气机\emph{曾经}所在的位置。
\end{solution}

\begin{solution}{exo:g7:sound-production:10}
水下的耳朵由水中 \qty{1500}{m/s} 的接力传递声音，
甲板上的耳朵则由空气中 \qty{340}{m/s} 的接力传递；说话者的大部分声音
根本不会穿过水面边界进入空气中。
介质越快，信息越隐蔽：尽管距离更远，
游泳者仍占优势。
\end{solution}

\begin{solution}{exo:g7:sound-production:11}
一次敲击，两条接力路径：经钢铁与经空气。钢铁中：
$3000 \div 5000 = \qty{0.6}{s}$；空气中：$3000 \div 340 \approx
\qty{8.8}{s}$——当啷声比轰响约早
\qty{8.2}{s}。
\end{solution}

\begin{solution}{exo:g7:sound-production:12}
一次敲击同时向两种介质发出信号；在
\qty{750}{m} 处，记录仪分别捕捉到水中到达的信号（水听器上）
和空气中到达的信号（麦克风上）。预期时间：
水中 $750 \div 1500 = \qty{0.5}{s}$，空气中 $750 \div 340
\approx \qty{2.2}{s}$——两者相差约 \qty{1.7}{s}。反过来推算：
已知 $d$ 和空气中的速率，测得的时间差
可得水中到达时间（$t_{\text{air}} -
\text{时间差}$），并由 $v = 750 \div 0.5 = \qty{1500}{m/s}$——
水中的速率就这样测得，从未对那次无声的敲击
本身计时。
\end{solution}

\begin{solution}{pb:g7:sound-production:1}
\textbf{1.} $340 \times 4 = \qty{1360}{m}$ 为往返路程，宽度
为 \qty{680}{m}。
\textbf{2.} $340 \times 1.5 = 510$；一半为 \qty{255}{m}。
\textbf{3.} 第一面墙：$340 \times 1 \div 2 = \qty{170}{m}$；
第二面墙：$340 \times 2 \div 2 = \qty{340}{m}$；峡谷沿该方向的跨度为
$170 + 340 = \qty{510}{m}$。
\textbf{4.} 公式是精确的，手却不精确：人为开始和停止计时的误差
可达十分之几秒，在 \qty{340}{m/s} 下相当于数十
米。因此有一条老办法：反复测量，看读数是否聚拢，
并报告这一聚拢的结果。
\textbf{5.} $1500 \times 0.2 = \qty{300}{m}$ 为往返路程，深度
为 \qty{150}{m}。
\textbf{6.} $1500 \times 0.4 \div 2 = \qty{300}{m}$。
\textbf{7.} 在空气与水的交界处，大部分呼喊声会被
反射回去——只有少量进入河水，而微弱的河底回声还必须
再次穿过交界才能到达空气中的耳朵：几乎
无法留存。即便承认有回声，往返时间只有十分之几秒，
也让任何用耳朵计时的方法无能为力。
\textbf{8.} 船驶过落差处时，回声时间会突然变长——
河底更深，往返时间更长：声呐
便能描绘出水下的悬崖。
\textbf{9.} 钢铁中：$2000 \div 5000 = \qty{0.4}{s}$；空气中：$2000
\div 340 \approx \qty{5.9}{s}$。
\textbf{10.} 两条接力同时出发，因此时间差与距离
成正比：每增加一千米，慢速接力就会固定多出一段
延迟。测出时间差，
除以每千米的时间差，距离便
迎刃而解——一次锤击即可，无需在远端计时。
\textbf{11.} 钢轨声：雾是由水滴组成的群体，它会
把\emph{光}彻底散射掉，却几乎不影响
固体钢铁内部的接力（空气中的声音在雾中也能较好地保留，
但钢轨传来的信息最清晰、最快，在三者中
居首）。
\textbf{12.} 例如：“每次回声都要为往返付出代价：
先减半再相信。每种介质都有自己的速率：
选速率须看介质本身，切勿凭习惯。这座峡谷中的每次计时
都始于眼睛——在此处，光的传播可视为
瞬间完成，因此看到锤子落下，便可充当
起跑信号。”
\end{solution}
